\centerline{\bf \Huge Слив региона ВсОШ}
\centerline{\bf \Huge 9 класс}

\begin{flushright}
        \scriptsize{Кто наберёт меньше 25 тот лох}
\end{flushright}

\problem{1}
В олимпиаде по математике было 10 задач. За каждую задачу Ваня получил хотя бы 1 балл. Он утверждает, что среднее арифметическое его баллов за любые несколько (хотя бы две) подряд идущих задач - нецелое число. Могли ли его слова оказаться правдой? Балл за каждую задачу - натуральное число от 1 до 7 .


\problem{2}
В компании некоторые пары людей являются друзьями (дружба взаимна). Известно, что если двое людей не являются друзьями, то они имеют ровно двух общих друзей. В этой компании нашлись двое друзей $A$ и $B$, у которых нет общих друзей. Докажите, что у $A$ и $B$ поровну друзей.

\problem{3}
Биссектриса угла $A B D$ пересекает основание $A D$ равнобокой трапеции $A B C D$ в точке $L$. Точки $K$ и $N$ на отрезках $A C$ и $C D$ выбраны соответственно так, что $A K=A L$ и $D N=D L$. Докажите, что точки $B, C, K, N$ лежат на одной окружности.

\problem{4}
Дан квадратный трёхчлен $f(x)$ и число $a$. Известно, что числа $a, f(a), f(f(a)), f(f(f(a)))$ образуют геометрическую прогрессию с положительным знаменателем именно в таком порядке. Докажите, что она постоянная.

\problem{5}
В таблице $20 \times 20$ отмечено 180 клеток таким образом, что никакие четыре из них не образуют квадрат $2 \times 2$. Докажите, что можно отметить ещё одну клетку с сохранением этого условия.